%%%%%%%% GRCON 2016 EXAMPLE LATEX SUBMISSION FILE %%%%%%%%%%%%%%%%%

\documentclass{article}

\usepackage{times}
\usepackage{graphicx} % more modern
\usepackage{subfigure} 
\usepackage{natbib}
\usepackage{hyperref}

% please use this for your draft submission
\usepackage[nohyperref]{grcon2016} 
% please use this for your camera ready paper
%\usepackage[accepted]{grcon2016}

\grcontitlerunning{Submission and Formatting Instructions for GRCON 2016}

\begin{document} 

\twocolumn[
\grcontitle{Submission and Formatting Instructions for \\
           GNU Radio Conference (GRCON 2016)}

\grconauthor{Your Name}{email@yourdomain.edu}
\grconaddress{Your Fantastic Institute,
            14142 Imaginary Way, Mountain View, CA 94043 USA}
\grconauthor{Your CoAuthor's Name}{email@coauthordomain.edu}
\grconaddress{Their Fantastic Institute,
            1571 Jay St., Arlington, VA 22203 USA}

\grconkeywords{software radio, gnu radio, dsp, GRCON}
]

\vskip 0.3in

\begin{abstract} 
The purpose of this document is to provide both the basic paper template and submission guidelines. 
Abstracts should succinct single paragraph, ideally between 4 to 6 sentences long.
\end{abstract} 

\section{Electronic Submission}
\label{submission}

Submission to GRCON 2016 will be entirely electronic.  
Information about the submission process and \LaTeX\ templates
are available on the conference web site at:

\begin{center}
\url{www.gnuradio.org/grcon-2016/papers/}
\end{center}

Questions and submission PDFs can be sent to 
\texttt{grcon@gnuradio.org}.

The guidelines below will be enforced for initial submissions and
camera-ready copies.  Here is a brief summary:
\begin{itemize}
\item Submissions must be in PDF.
\item The maximum paper length is \textbf{8 pages excluding references and acknowledgements, and 10 pages
  including references and acknowledgements} (pages 9 and 10 must contain only references and acknowledgements).

\item Your paper should be in \textbf{10 point Times font}.
\item Place figure captions {\em under} the figure (and omit titles from inside graphics).  Place table captions {\em over} the table.
\item Do not alter the style template; in particular, do not compress the paper format by reducing the vertical spaces.
\item Keep your abstract brief and self-contained, one paragraph and roughly 4-6 sentences.  Title should have content words capitalized.
\item Where possible, provide source code and example implementation via reference to your github account.
\end{itemize}

\subsection{Submitting Papers}

{\bf Paper Deadline:} The deadline for paper submission to GRCON 2016
is \textbf{July 1, 2016}. 

{\bf Simultaneous Submission:} GRCON papers must not have been previously 
published or under review elsewhere.  Pre-publication to Arxiv is encouraged.

\medskip

Authors using \textbf{Word} must convert their document to PDF.  
Submissions will only be accepted in PDF format.

\subsection{Feedback from Paper Submission}

Acceptance notification will occur within one week of paper submission via email.
Unless otherwise specified, conference abstracts must previously be accepted.

\subsection{Submitting Final Camera-Ready Copy}

The final versions of papers accepted for publication should follow the
same format as initial submissions, but should use the accepted keyword 
for the style package.

Simply change:
$\mathtt{\backslash usepackage\{grcon2016\}}$ to 

$$\mathtt{\backslash usepackage[accepted]\{grcon2016\}}$$

\noindent

Camera-ready copies should have the title of the paper as running head
on each page except the first one.  

\section{Format of the Paper} 
 
All submissions must follow the same format to ensure consistent reproduction
and to let readers more easily find information.

\subsection{Length and Dimensions}

Papers must not exceed eight (8) pages of content, and ten (10) pages including 
acknowledgements and references.

The text of the paper should be formatted in two columns, with an
overall width of 6.75 inches, height of 9.0 inches, and 0.25 inches
between the columns. The left margin should be 0.75 inches and the top
margin 1.0 inch (2.54~cm) and US Letter size.

The paper body should be set in 10~point type with a vertical spacing
of 11~points. Please use Times typeface throughout the text.

\subsection{Title}

The paper title should be set in 14~point bold type and centered
between two horizontal rules that are 1~point thick, with 1.0~inch
between the top rule and the top edge of the page. Capitalize the
first letter of content words and put the rest of the title in lower
case.

\subsection{Author Information for Submission}
\label{author info}

Author information should include full names, institutional affiliation,
and at least city address, as well as email or web addresses.

\subsection{Abstract}

The paper abstract should begin in the left column, 0.4~inches below
the final address. The heading `Abstract' should be centered, bold,
and in 11~point type. The abstract body should use 10~point type, with
a vertical spacing of 11~points, and should be indented 0.25~inches
more than normal on left-hand and right-hand margins. 
\textbf{Keep your abstract brief and self-contained.}

\subsection{Partitioning the Text} 

You should organize your paper into sections and paragraphs to help
readers place a structure on the material and understand its
contributions.

\subsubsection{Sections and Subsections}

Section headings should be numbered, flush left, and set in 11~pt bold
type with the content words capitalized. Leave 0.25~inches of space
before the heading.

Similarly, subsection headings should be numbered, flush left, and set
in 10~pt bold type with the content words capitalized. Leave
0.2~inches of space before the heading.

Subsubsection headings should be numbered, flush left, and
set in 10~pt small caps with the content words capitalized. Leave
0.18~inches of space before the heading.

Please use no more than three levels of headings.

\subsubsection{Paragraphs and Footnotes}

Within each section or subsection, you should further partition the
paper into paragraphs. Do not indent the first line of a given
paragraph, but insert a blank line between succeeding ones.

\begin{figure}[ht]
  \begin{center}
    \centerline{\includegraphics[width=\columnwidth]{gr.png}}
    \caption{GNU Radio Logo}
    \label{grlogo}
  \end{center}
\end{figure} 

\subsection{Figures}
 
You may want to include figures in the paper to help readers visualize
your approach and your results. Such artwork should be centered,
legible, and clearly labeled and explained with captions similar to Figure~\ref{grlogo}.

\subsection{Citations and References} 

Please use APA reference format regardless of your formatter
or word processor. If you rely on the \LaTeX\/ bibliographic 
facility, use {\tt natbib.sty} and {\tt grcon2016.bst} 
included in the style-file package to obtain this format.

This template is based off the excellent template for ICML provided
in the work by \cite{langley00}.

\subsection{Software and Data}

We strongly encourage the publication of software and data with the
camera-ready version of the paper whenever appropriate.  This can be
done by including a URL in the camera-ready copy.  

\bibliography{example_paper}
\bibliographystyle{grcon2016}

\end{document} 
